\documentclass[12pt,a4paper]{article}
\usepackage[UTF8]{ctex}
\usepackage{amsmath,amssymb,amsthm}
\usepackage{geometry}
\usepackage{graphicx}
\usepackage{tikz}
\usepackage{pgfplots}
\usetikzlibrary{shapes,arrows,positioning,calc}
\geometry{left=2.5cm,right=2.5cm,top=2.5cm,bottom=2.5cm}

\title{12.1.7.3 测量形状闭合抓取质量详解}
\author{}
\date{}

\begin{document}

\maketitle

\section{引言}

\subsection{问题提出}

考虑图12.16中显示的两个形状闭合抓取。\textbf{哪个是更好的抓取？}

\begin{itemize}
\item 两个抓取都实现了形状闭合
\item 但它们的\textbf{质量}可能不同
\item 需要一种\textbf{度量}来量化抓取质量
\end{itemize}

\subsection{抓取度量的定义}

\textbf{抓取度量}：取接触集$\{F_i\}$，返回单个值$\text{Qual}(\{F_i\})$

\begin{itemize}
\item $\text{Qual}(\{F_i\}) < 0$：抓取\textbf{不产生}形状闭合
\item $\text{Qual}(\{F_i\}) = 0$：抓取\textbf{刚好}产生形状闭合（临界情况）
\item $\text{Qual}(\{F_i\}) > 0$：抓取产生形状闭合，\textbf{值越大越好}
\end{itemize}

\section{接触Wrench集合 $\mathcal{C}_F$}

\subsection{物理约束}

\textbf{假设}：为了避免损坏物体，接触$i$处的力大小必须满足：
\begin{equation}
|f_i| \leq f_{i,\max} > 0
\label{eq:force_limit}
\end{equation}

\textbf{物理解释}：
\begin{itemize}
\item 每个接触有\textbf{最大力限制}
\item 这防止物体被损坏
\item 也防止手指施加过大的力
\end{itemize}

\subsection{接触Wrench集合的定义}

\textbf{定义}：$j$个接触可以施加的接触wrench集合为：
\begin{equation}
\mathcal{C}_F = \left\{\sum_{i=1}^j f_i F_i \middle| f_i \in [0, f_{i,\max}] \right\}
\label{eq:12.15}
\end{equation}

\textbf{关键点}：
\begin{itemize}
\item $f_i \in [0, f_{i,\max}]$：力必须\textbf{非负}（只能推，不能拉）
\item $\sum_{i=1}^j f_i F_i$：所有接触wrench的\textbf{正线性组合}
\item $\mathcal{C}_F$是\textbf{凸多面体}（polytope）
\end{itemize}

\subsection{几何解释}

\textbf{图12.17的说明}：

\begin{itemize}
\item \textbf{二维情况}：$\mathcal{C}_F$是wrench空间中的\textbf{凸多边形}
\item \textbf{顶点}：由某些$f_i = f_{i,\max}$或$f_i = 0$的组合产生
\item \textbf{边}：由约束$f_i \in [0, f_{i,\max}]$的边界产生
\end{itemize}

\textbf{形状闭合的几何特征}：

如果抓取产生形状闭合，则：
\begin{itemize}
\item $\mathcal{C}_F$在其\textbf{内部}包含wrench空间的原点
\item 这意味着：接触可以产生\textbf{任何方向}的wrench来抵抗干扰
\end{itemize}

\section{从多面体到质量值}

\subsection{问题}

\textbf{挑战}：如何将凸多面体$\mathcal{C}_F$转换为单个数字$\text{Qual}(\{F_i\})$？

\textbf{理想方法}：
\begin{itemize}
\item 考虑物体可能遇到的\textbf{干扰wrench分布}
\item 基于实际应用场景设计度量
\end{itemize}

\textbf{简化方法}：
\begin{itemize}
\item 使用\textbf{最大内接球}方法
\item 计算以原点为中心、适合$\mathcal{C}_F$内部的最大球的半径
\end{itemize}

\subsection{最大内接球方法}

\textbf{定义}：
\begin{equation}
\text{Qual}(\{F_i\}) = \text{最大内接球的半径}
\end{equation}

\textbf{几何意义}（图12.17）：
\begin{itemize}
\item 以wrench空间原点为中心
\item 完全位于$\mathcal{C}_F$内部
\item 半径越大，抓取质量越好
\end{itemize}

\textbf{物理解释}：
\begin{itemize}
\item 半径$d$表示：接触可以抵抗\textbf{任何方向}、大小不超过$d$的干扰wrench
\item 更大的半径意味着更强的\textbf{鲁棒性}
\end{itemize}

\section{两个关键问题}

\subsection{问题1：单位不一致}

\textbf{问题}：力矩和力具有\textbf{不同的单位}

\begin{itemize}
\item 力：牛顿（N）
\item 力矩：牛顿·米（N·m）
\item 如何比较它们的大小？
\end{itemize}

\textbf{解决方案}：引入\textbf{特征长度}$r$

\begin{itemize}
\item 选择被抓取物体的特征长度$r$（例如，物体尺寸）
\item 将接触力矩$m$转换为等效力：$m/r$
\item 这样力矩和力就有了\textbf{相同的单位}
\end{itemize}

\textbf{数学表达}：

对于平面情况，wrench原本是：
\begin{equation}
F_i = \begin{bmatrix} m_{iz} \\ f_{ix} \\ f_{iy} \end{bmatrix}
\end{equation}

转换后：
\begin{equation}
\tilde{F}_i = \begin{bmatrix} m_{iz}/r \\ f_{ix} \\ f_{iy} \end{bmatrix}
\end{equation}

现在所有分量都有力的单位。

\subsection{问题2：坐标系原点依赖}

\textbf{问题}：接触力产生的力矩\textbf{依赖于}空间坐标系原点的位置

\begin{itemize}
\item 如果原点改变，力矩分量会改变
\item 这会影响$\mathcal{C}_F$的形状
\item 进而影响质量度量
\end{itemize}

\textbf{解决方案}：选择合适的原点位置

\textbf{常见选择}：
\begin{itemize}
\item \textbf{几何中心}：物体形状的中心
\item \textbf{质心}：物体的质量中心
\item \textbf{抓取中心}：接触点的中心
\end{itemize}

\textbf{建议}：
\begin{itemize}
\item 选择在物体\textbf{几何中心附近}
\item 或选择在物体\textbf{质心}处
\item 这样度量更稳定、更有意义
\end{itemize}

\section{边界超平面详解}

\subsection{什么是超平面？}

\textbf{超平面}是$n$维空间中的$(n-1)$维子空间。

\textbf{从低维到高维的理解}：
\begin{itemize}
\item \textbf{一维空间}（直线）：超平面是\textbf{点}（0维）
\item \textbf{二维空间}（平面）：超平面是\textbf{直线}（1维）
\item \textbf{三维空间}：超平面是\textbf{平面}（2维）
\item \textbf{$n$维空间}：超平面是$(n-1)$维子空间
\end{itemize}

\textbf{超平面的方程}：

在$n$维空间中，超平面可以表示为：
\begin{equation}
a_1 x_1 + a_2 x_2 + \cdots + a_n x_n = b
\label{eq:hyperplane_general}
\end{equation}

或者用向量形式：
\begin{equation}
a^T x = b
\label{eq:hyperplane_vector}
\end{equation}

其中：
\begin{itemize}
\item $a = [a_1, a_2, \ldots, a_n]^T$是超平面的\textbf{法向量}
\item $b$是标量常数
\item $x = [x_1, x_2, \ldots, x_n]^T$是空间中的点
\end{itemize}

\subsection{具体例子}

\textbf{例子1：二维空间中的超平面（直线）}

在二维空间中，超平面是直线：
\begin{equation}
a_1 x_1 + a_2 x_2 = b
\end{equation}

例如：$x_1 + x_2 = 1$是一条直线。

\textbf{例子2：三维空间中的超平面（平面）}

在三维空间中，超平面是平面：
\begin{equation}
a_1 x_1 + a_2 x_2 + a_3 x_3 = b
\end{equation}

例如：$x_1 + x_2 + x_3 = 1$是一个平面。

\textbf{例子3：wrench空间中的超平面}

对于平面物体的wrench空间（3维），超平面是2维平面：
\begin{equation}
a_1 m_z + a_2 f_x + a_3 f_y = b
\end{equation}

\subsection{什么是边界超平面？}

\textbf{边界超平面}是多面体（polytope）的\textbf{边界}。

\textbf{多面体}：
\begin{itemize}
\item 由\textbf{半空间}的交集定义
\item 半空间：$a^T x \leq b$或$a^T x \geq b$
\item 边界：满足$a^T x = b$的点
\end{itemize}

\textbf{边界超平面的定义}：

如果多面体$\mathcal{C}_F$由以下约束定义：
\begin{equation}
\mathcal{C}_F = \{x \in \mathbb{R}^n \mid a_i^T x \leq b_i, i = 1, \ldots, m\}
\end{equation}

那么每个约束$a_i^T x = b_i$对应一个\textbf{边界超平面}。

\subsection{$\mathcal{C}_F$的边界超平面}

\textbf{回顾}：接触wrench集合定义为：
\begin{equation}
\mathcal{C}_F = \left\{\sum_{i=1}^j f_i F_i \middle| f_i \in [0, f_{i,\max}] \right\}
\end{equation}

\textbf{约束条件}：
\begin{itemize}
\item $f_i \geq 0$（力不能为负）
\item $f_i \leq f_{i,\max}$（力有上限）
\end{itemize}

\textbf{边界超平面的来源}：

每个约束的\textbf{边界}对应一个超平面：
\begin{itemize}
\item $f_i = 0$：第$i$个接触的力为零
\item $f_i = f_{i,\max}$：第$i$个接触的力达到最大值
\end{itemize}

\textbf{在wrench空间中的表示}：

设$F = [F_1, F_2, \ldots, F_j]$，则：
\begin{equation}
x = \sum_{i=1}^j f_i F_i = F f
\end{equation}

其中$f = [f_1, f_2, \ldots, f_j]^T$。

约束$f_i = 0$在wrench空间中对应：
\begin{equation}
x = \sum_{k \neq i} f_k F_k
\end{equation}

这是一个\textbf{超平面}（由其他接触wrench张成的子空间）。

约束$f_i = f_{i,\max}$在wrench空间中对应：
\begin{equation}
x = f_{i,\max} F_i + \sum_{k \neq i} f_k F_k
\end{equation}

这也是一个\textbf{超平面}（平移后的子空间）。

\subsection{二维例子：多边形}

\textbf{情况}：二维wrench空间，三个接触

\begin{itemize}
\item $\mathcal{C}_F$是一个\textbf{凸多边形}
\item 边界由\textbf{边}（1维超平面，即直线）组成
\item 每条边对应某个约束的边界
\end{itemize}

\textbf{具体例子}：

假设$f_1, f_2, f_3 \in [0, 1]$，则：
\begin{itemize}
\item 边1：$f_1 = 0$，$f_2, f_3 \in [0,1]$
\item 边2：$f_1 = 1$，$f_2, f_3 \in [0,1]$
\item 边3：$f_2 = 0$，$f_1, f_3 \in [0,1]$
\item 边4：$f_2 = 1$，$f_1, f_3 \in [0,1]$
\item 边5：$f_3 = 0$，$f_1, f_2 \in [0,1]$
\item 边6：$f_3 = 1$，$f_1, f_2 \in [0,1]$
\end{itemize}

\textbf{注意}：不是所有约束边界都会出现在多边形的边上（有些可能在内部）。

\subsection{超平面的法向量}

\textbf{法向量}：垂直于超平面的向量

对于超平面$a^T x = b$：
\begin{itemize}
\item 法向量：$a$（或$-a$）
\item 如果$|a| = 1$，则$b$是原点到超平面的\textbf{有符号距离}
\end{itemize}

\textbf{符号约定}：
\begin{itemize}
\item 如果$b > 0$，原点在超平面的\textbf{正侧}（法向量指向原点）
\item 如果$b < 0$，原点在超平面的\textbf{负侧}（法向量背离原点）
\item 如果$b = 0$，原点在超平面上
\end{itemize}

\subsection{如何找到边界超平面？}

\textbf{方法1：从约束条件}

\begin{enumerate}
\item 列出所有约束：$f_i \geq 0$和$f_i \leq f_{i,\max}$
\item 对每个约束，设等号成立：$f_i = 0$或$f_i = f_{i,\max}$
\item 在wrench空间中，这些约束对应超平面
\end{enumerate}

\textbf{方法2：从多面体表示}

如果$\mathcal{C}_F = \{x \mid A x \leq b\}$，则：
\begin{itemize}
\item 每行$A_i$对应一个约束
\item 边界超平面：$A_i^T x = b_i$
\end{itemize}

\textbf{方法3：从顶点表示}

\begin{itemize}
\item 找到多面体的所有顶点
\item 找到连接顶点的边/面
\item 这些边/面就是边界超平面
\end{itemize}

\section{质量度量的计算}

\subsection{有符号距离}

\textbf{给定}：
\begin{itemize}
\item 空间坐标系的选择（原点位置）
\item 特征长度$r$的选择
\item 接触wrench集合$\mathcal{C}_F$
\end{itemize}

\textbf{方法}：计算从wrench空间原点到$\mathcal{C}_F$边界上每个\textbf{超平面}的有符号距离

\subsection{超平面方程}

\textbf{多面体$\mathcal{C}_F$的边界}由超平面组成。

每个超平面可以表示为：
\begin{equation}
a^T x = b
\end{equation}

其中：
\begin{itemize}
\item $a$是超平面的法向量（单位向量）
\item $b$是超平面到原点的距离
\item $x$是wrench空间中的点
\end{itemize}

\subsection{点到超平面的距离}

\textbf{有符号距离}：
\begin{equation}
d = \frac{b}{|a|} = b \quad \text{（如果$|a| = 1$）}
\end{equation}

\textbf{符号意义}：
\begin{itemize}
\item $d > 0$：原点在超平面的\textbf{正侧}（法向量指向原点）
\item $d < 0$：原点在超平面的\textbf{负侧}（法向量背离原点）
\item $d = 0$：原点在超平面上
\end{itemize}

\subsection{质量度量的定义}

\textbf{定义}：
\begin{equation}
\text{Qual}(\{F_i\}) = \min\{d_1, d_2, \ldots, d_k\}
\label{eq:quality_metric}
\end{equation}

其中$d_1, d_2, \ldots, d_k$是从原点到$\mathcal{C}_F$边界上所有超平面的有符号距离。

\textbf{关键理解}：
\begin{itemize}
\item 取\textbf{最小值}确保球完全在$\mathcal{C}_F$内部
\item 如果所有$d_i > 0$，则原点在$\mathcal{C}_F$内部（形状闭合）
\item 如果某个$d_i < 0$，则原点不在$\mathcal{C}_F$内部（不是形状闭合）
\end{itemize}

\section{图12.17的详细解释}

\subsection{图12.17(a)}

\textbf{情况}：三个接触wrench在二维wrench空间中

\begin{itemize}
\item $\mathcal{C}_F$是一个\textbf{凸多边形}
\item 以原点为中心的最大内接球半径为$d$
\item $d$是到所有边的最小距离
\end{itemize}

\textbf{质量}：$\text{Qual}(\{F_i\}) = d$

\subsection{图12.17(b)}

\textbf{情况}：不同的三个接触wrench

\begin{itemize}
\item $\mathcal{C}_F$的形状不同
\item 最大内接球的半径\textbf{更大}
\item 说明这个抓取质量\textbf{更好}
\end{itemize}

\textbf{质量}：$\text{Qual}(\{F_i\}) = d' > d$

\section{图12.16的例子}

\subsection{两个抓取的比较}

\textbf{假设}：每个手指允许施加相同的力（$f_{i,\max}$相同）

\textbf{观察}：
\begin{itemize}
\item 左侧的抓取：接触可以抵抗关于物体中心的\textbf{更大力矩}
\item 右侧的抓取：抵抗力矩的能力较弱
\end{itemize}

\textbf{结论}：
\begin{itemize}
\item 左侧抓取的$\mathcal{C}_F$更大
\item 最大内接球半径更大
\item 因此\textbf{质量更好}
\end{itemize}

\section{算法步骤总结}

\subsection{计算抓取质量的步骤}

\textbf{步骤1}：选择坐标系和特征长度
\begin{itemize}
\item 选择空间坐标系原点（几何中心或质心）
\item 选择特征长度$r$（物体尺寸）
\end{itemize}

\textbf{步骤2}：构建接触wrench集合
\begin{equation}
\mathcal{C}_F = \left\{\sum_{i=1}^j f_i F_i \middle| f_i \in [0, f_{i,\max}] \right\}
\end{equation}

\textbf{步骤3}：单位归一化
\begin{itemize}
\item 将力矩分量除以$r$：$m_{iz} \to m_{iz}/r$
\item 确保所有分量单位一致
\end{itemize}

\textbf{步骤4}：找到$\mathcal{C}_F$的边界超平面
\begin{itemize}
\item 确定多面体的所有面（超平面）
\item 每个面由约束$f_i = 0$或$f_i = f_{i,\max}$定义
\end{itemize}

\textbf{步骤5}：计算有符号距离
\begin{itemize}
\item 对每个超平面，计算原点到它的有符号距离
\item $d_i = \text{有符号距离到第$i$个超平面}$
\end{itemize}

\textbf{步骤6}：计算质量度量
\begin{equation}
\text{Qual}(\{F_i\}) = \min_i d_i
\end{equation}

\section{实际应用考虑}

\subsection{度量的选择}

\textbf{注意}：最大内接球方法只是\textbf{一种}选择

\textbf{其他可能的度量}：
\begin{itemize}
\item \textbf{体积}：$\mathcal{C}_F$的体积
\item \textbf{最小距离}：到最近边界的距离（就是我们的方法）
\item \textbf{条件数}：基于$F$矩阵的条件数
\item \textbf{力传递比}：基于力放大系数
\end{itemize}

\subsection{度量的优缺点}

\textbf{最大内接球方法的优点}：
\begin{itemize}
\item \textbf{几何直观}：容易理解和可视化
\item \textbf{方向无关}：对所有方向的干扰都公平
\item \textbf{计算简单}：只需要计算距离
\end{itemize}

\textbf{最大内接球方法的缺点}：
\begin{itemize}
\item \textbf{不考虑干扰分布}：假设所有方向等概率
\item \textbf{可能过于保守}：实际干扰可能只在某些方向
\item \textbf{单位选择敏感}：特征长度$r$的选择影响结果
\end{itemize}

\section{数学细节补充}

\subsection{多面体的表示}

\textbf{$\mathcal{C}_F$的顶点表示}：

$\mathcal{C}_F$的顶点由以下方式产生：
\begin{itemize}
\item 对每个接触$i$，选择$f_i = 0$或$f_i = f_{i,\max}$
\item 共有$2^j$种组合（但很多可能不在边界上）
\end{itemize}

\textbf{$\mathcal{C}_F$的半空间表示}：

$\mathcal{C}_F$可以表示为：
\begin{equation}
\mathcal{C}_F = \{x \in \mathbb{R}^n \mid A x \leq b\}
\end{equation}

其中：
\begin{itemize}
\item $A$是约束矩阵
\item $b$是约束向量
\item 每行对应一个约束（$f_i \geq 0$或$f_i \leq f_{i,\max}$）
\end{itemize}

\subsection{有符号距离的计算}

\textbf{对于半空间表示}：

如果$\mathcal{C}_F = \{x \mid a^T x \leq b\}$，则：
\begin{itemize}
\item 超平面方程：$a^T x = b$
\item 原点到超平面的距离：$d = \frac{b}{|a|}$
\item 如果$|a| = 1$（归一化），则$d = b$
\end{itemize}

\textbf{符号判断}：
\begin{itemize}
\item 如果$a^T 0 = 0 \leq b$，则原点在$\mathcal{C}_F$内，$d = b > 0$
\item 如果$a^T 0 = 0 > b$，则原点不在$\mathcal{C}_F$内，$d = b < 0$
\end{itemize}

\section{具体计算例子：二维平面抓取}

\subsection{问题设置}

\textbf{场景}：在二维平面中抓取一个矩形物体

\textbf{接触配置}：
\begin{itemize}
\item 4个接触点（满足形状闭合的最小数量）
\item 接触位置：矩形的四个角附近
\item 每个接触的法向力指向物体内部
\end{itemize}

\textbf{具体参数}：

假设矩形物体尺寸为$2 \times 1$（单位：米），质心在原点。

四个接触点的位置和法向量：
\begin{itemize}
\item 接触1：位置$(1, 0.5)$，法向量$\hat{n}_1 = (-1, 0)$（向右推）
\item 接触2：位置$(-1, 0.5)$，法向量$\hat{n}_2 = (1, 0)$（向左推）
\item 接触3：位置$(1, -0.5)$，法向量$\hat{n}_3 = (-1, 0)$（向右推）
\item 接触4：位置$(-1, -0.5)$，法向量$\hat{n}_4 = (1, 0)$（向左推）
\end{itemize}

\textbf{注意}：这是一个简化的例子，实际抓取中接触法向量可能不完全是水平方向。

\subsection{计算接触Wrench}

\textbf{对于平面情况}，wrench为：
\begin{equation}
F_i = \begin{bmatrix} m_{iz} \\ f_{ix} \\ f_{iy} \end{bmatrix}
\end{equation}

其中：
\begin{itemize}
\item $m_{iz}$：关于$z$轴的力矩
\item $f_{ix}, f_{iy}$：力的$x$和$y$分量
\end{itemize}

\textbf{接触wrench的计算}：

对于点接触，wrench为：
\begin{equation}
F_i = \begin{bmatrix} 
(p_i \times \hat{n}_i)_z \\
\hat{n}_{ix} \\
\hat{n}_{iy}
\end{bmatrix}
\end{equation}

其中$p_i$是接触点位置。

\textbf{具体计算}：

\textbf{接触1}：$p_1 = (1, 0.5)$，$\hat{n}_1 = (-1, 0)$
\begin{align}
m_{1z} &= (p_1 \times \hat{n}_1)_z = 1 \cdot 0 - 0.5 \cdot (-1) = 0.5 \\
F_1 &= \begin{bmatrix} 0.5 \\ -1 \\ 0 \end{bmatrix}
\end{align}

\textbf{接触2}：$p_2 = (-1, 0.5)$，$\hat{n}_2 = (1, 0)$
\begin{align}
m_{2z} &= (p_2 \times \hat{n}_2)_z = (-1) \cdot 0 - 0.5 \cdot 1 = -0.5 \\
F_2 &= \begin{bmatrix} -0.5 \\ 1 \\ 0 \end{bmatrix}
\end{align}

\textbf{接触3}：$p_3 = (1, -0.5)$，$\hat{n}_3 = (-1, 0)$
\begin{align}
m_{3z} &= (p_3 \times \hat{n}_3)_z = 1 \cdot 0 - (-0.5) \cdot (-1) = -0.5 \\
F_3 &= \begin{bmatrix} -0.5 \\ -1 \\ 0 \end{bmatrix}
\end{align}

\textbf{接触4}：$p_4 = (-1, -0.5)$，$\hat{n}_4 = (1, 0)$
\begin{align}
m_{4z} &= (p_4 \times \hat{n}_4)_z = (-1) \cdot 0 - (-0.5) \cdot 1 = 0.5 \\
F_4 &= \begin{bmatrix} 0.5 \\ 1 \\ 0 \end{bmatrix}
\end{align}

\subsection{构建接触Wrench集合}

\textbf{假设}：每个接触的最大力$f_{i,\max} = 1$（单位：N）

\textbf{接触wrench集合}：
\begin{equation}
\mathcal{C}_F = \left\{\sum_{i=1}^4 f_i F_i \middle| f_i \in [0, 1] \right\}
\end{equation}

\textbf{wrench矩阵}：
\begin{equation}
F = \begin{bmatrix}
F_1 & F_2 & F_3 & F_4
\end{bmatrix} = \begin{bmatrix}
0.5 & -0.5 & -0.5 & 0.5 \\
-1 & 1 & -1 & 1 \\
0 & 0 & 0 & 0
\end{bmatrix}
\end{equation}

\textbf{注意}：这个例子中所有接触的$y$方向力分量为0，所以wrench空间实际上是2维的（力矩和$x$方向力）。

\subsection{简化：二维Wrench空间}

由于$f_{iy} = 0$对所有接触，我们可以将问题简化为2维：
\begin{equation}
\tilde{F}_i = \begin{bmatrix} m_{iz} \\ f_{ix} \end{bmatrix}
\end{equation}

\textbf{简化的wrench矩阵}：
\begin{equation}
\tilde{F} = \begin{bmatrix}
0.5 & -0.5 & -0.5 & 0.5 \\
-1 & 1 & -1 & 1
\end{bmatrix}
\end{equation}

\textbf{简化的接触wrench集合}：
\begin{equation}
\mathcal{C}_F = \left\{\sum_{i=1}^4 f_i \tilde{F}_i \middle| f_i \in [0, 1] \right\}
\end{equation}

\subsection{找到边界超平面（直线）}

在2维空间中，超平面是\textbf{直线}。

\textbf{方法}：找到$\mathcal{C}_F$的顶点

$\mathcal{C}_F$的顶点由$f_i \in \{0, 1\}$的组合产生。

\textbf{计算所有顶点}：

设$f = [f_1, f_2, f_3, f_4]^T$，则$w = \tilde{F} f$。

\textbf{关键顶点}（选择几个重要的）：
\begin{itemize}
\item $f = [1, 0, 0, 0]^T$：$w_1 = [0.5, -1]^T$
\item $f = [0, 1, 0, 0]^T$：$w_2 = [-0.5, 1]^T$
\item $f = [0, 0, 1, 0]^T$：$w_3 = [-0.5, -1]^T$
\item $f = [0, 0, 0, 1]^T$：$w_4 = [0.5, 1]^T$
\item $f = [1, 1, 0, 0]^T$：$w_5 = [0, 0]^T$（原点！）
\item $f = [1, 0, 1, 0]^T$：$w_6 = [0, -2]^T$
\item $f = [0, 1, 0, 1]^T$：$w_7 = [0, 2]^T$
\item $f = [1, 0, 0, 1]^T$：$w_8 = [1, 0]^T$
\item $f = [0, 1, 1, 0]^T$：$w_9 = [-1, 0]^T$
\end{itemize}

\textbf{观察}：$w_5 = [0, 0]^T$是原点，说明这个抓取\textbf{可以产生形状闭合}（因为原点在$\mathcal{C}_F$内部）。

\subsection{确定多边形的边界}

\textbf{凸包}：$\mathcal{C}_F$是这些顶点的凸包。

\textbf{关键观察}：
\begin{itemize}
\item 原点$[0, 0]^T$在$\mathcal{C}_F$内部
\item 边界由连接顶点的直线段组成
\end{itemize}

\textbf{边界直线}（简化分析）：

考虑约束$f_i = 0$和$f_i = 1$对应的边界：

\textbf{边界1}：$f_1 = 0$，其他$f_i \in [0,1]$
\begin{equation}
w = f_2 \begin{bmatrix} -0.5 \\ 1 \end{bmatrix} + f_3 \begin{bmatrix} -0.5 \\ -1 \end{bmatrix} + f_4 \begin{bmatrix} 0.5 \\ 1 \end{bmatrix}
\end{equation}

这是一个\textbf{三角形区域}的边界。

\textbf{更简单的方法}：直接计算原点到各约束边界的距离。

\subsection{单位归一化}

\textbf{特征长度}：选择$r = 1$米（物体的特征尺寸）

\textbf{归一化wrench}：
\begin{equation}
\tilde{F}_i^{\text{norm}} = \begin{bmatrix} m_{iz}/r \\ f_{ix} \end{bmatrix} = \begin{bmatrix} m_{iz} \\ f_{ix} \end{bmatrix}
\end{equation}

（因为$r = 1$，所以不变）

\subsection{计算有符号距离}

\textbf{方法}：对于每个约束边界，计算原点到对应超平面的距离。

\textbf{约束边界对应的超平面}：

在2维wrench空间中，我们需要找到$\mathcal{C}_F$的\textbf{支撑超平面}（supporting hyperplanes）。

\textbf{简化方法}：使用线性规划找到最近边界

\textbf{问题}：找到最小的$d$，使得以原点为中心、半径为$d$的球完全在$\mathcal{C}_F$内部。

这等价于：
\begin{align}
\text{最大化} \quad & d \\
\text{使得} \quad & \|w\| \leq d \text{ 的所有 } w \in \mathcal{C}_F \text{ 都满足约束}
\end{align}

\textbf{更实用的方法}：检查原点是否在$\mathcal{C}_F$内部，并计算到最近边界的距离。

\subsection{使用半空间表示}

\textbf{$\mathcal{C}_F$的半空间表示}：

$\mathcal{C}_F$可以表示为：
\begin{equation}
\mathcal{C}_F = \{w \in \mathbb{R}^2 \mid \exists f \in [0,1]^4, w = \tilde{F} f\}
\end{equation}

\textbf{对偶表示}：找到所有支撑超平面。

\textbf{关键支撑超平面}：

考虑从原点出发的各个方向，找到$\mathcal{C}_F$在这些方向上的\textbf{支撑点}。

\textbf{方向} $\theta$对应的支撑点：
\begin{equation}
w^*(\theta) = \arg\max_{w \in \mathcal{C}_F} \cos(\theta) w_1 + \sin(\theta) w_2
\end{equation}

\textbf{支撑超平面}：垂直于方向$\theta$，通过$w^*(\theta)$。

\textbf{原点到支撑超平面的距离}：
\begin{equation}
d(\theta) = \cos(\theta) w_1^*(\theta) + \sin(\theta) w_2^*(\theta)
\end{equation}

\subsection{数值计算示例}

\textbf{关键方向}：

\textbf{方向1}：$\theta = 0$（$x$轴正方向）
\begin{itemize}
\item 最大化$w_1$（力矩分量）
\item 约束：$w = \tilde{F} f$，$f_i \in [0,1]$
\item 解：$f = [1, 0, 0, 1]^T$，$w^* = [1, 0]^T$
\item 距离：$d(0) = 1$
\end{itemize}

\textbf{方向2}：$\theta = \pi$（$x$轴负方向）
\begin{itemize}
\item 最大化$-w_1$
\item 解：$f = [0, 1, 1, 0]^T$，$w^* = [-1, 0]^T$
\item 距离：$d(\pi) = 1$
\end{itemize}

\textbf{方向3}：$\theta = \pi/2$（$y$轴正方向，力$x$分量正方向）
\begin{itemize}
\item 最大化$w_2$（力$x$分量）
\item 解：$f = [0, 1, 0, 1]^T$，$w^* = [0, 2]^T$
\item 距离：$d(\pi/2) = 2$
\end{itemize}

\textbf{方向4}：$\theta = -\pi/2$（$y$轴负方向，力$x$分量负方向）
\begin{itemize}
\item 最大化$-w_2$
\item 解：$f = [1, 0, 1, 0]^T$，$w^* = [0, -2]^T$
\item 距离：$d(-\pi/2) = 2$
\end{itemize}

\textbf{方向5}：$\theta = \pi/4$
\begin{itemize}
\item 最大化$(w_1 + w_2)/\sqrt{2}$
\item 需要求解线性规划
\item 近似：考虑顶点，$f = [1, 0, 0, 1]^T$给出$w = [1, 0]^T$
\item 但可能不是最优，需要检查其他组合
\end{itemize}

\textbf{关键观察}：最小距离可能不在这些主方向上。

\subsection{简化计算：使用对称性}

\textbf{对称性分析}：

由于抓取配置的对称性，$\mathcal{C}_F$关于原点对称（实际上，原点在内部）。

\textbf{关键边界}：

考虑约束$f_1 + f_2 = 1$和$f_3 + f_4 = 1$的情况（内力平衡）。

\textbf{特殊情况}：$f_1 = f_2 = 0.5$，$f_3 = f_4 = 0.5$
\begin{align}
w &= 0.5 \begin{bmatrix} 0.5 \\ -1 \end{bmatrix} + 0.5 \begin{bmatrix} -0.5 \\ 1 \end{bmatrix} + 0.5 \begin{bmatrix} -0.5 \\ -1 \end{bmatrix} + 0.5 \begin{bmatrix} 0.5 \\ 1 \end{bmatrix} \\
&= \begin{bmatrix} 0 \\ 0 \end{bmatrix}
\end{align}

这确认了原点在$\mathcal{C}_F$内部。

\subsection{系统化计算质量度量}

\textbf{方法}：使用线性规划找到最小距离

\textbf{问题表述}：

找到最大的$d$，使得以原点为中心、半径为$d$的球完全在$\mathcal{C}_F$内部。

等价于：
\begin{align}
\text{最大化} \quad & d \\
\text{使得} \quad & \text{对于所有方向 } \theta, \text{ 有 } d \leq d(\theta)
\end{align}

其中$d(\theta)$是方向$\theta$上的支撑距离。

\textbf{关键方向检查}：

\textbf{方向1}：$\theta = 0$（正力矩方向）
\begin{itemize}
\item 线性规划：$\max w_1$，约束$w = \tilde{F} f$，$f_i \in [0,1]$
\item 解：$f = [1, 0, 0, 1]^T$
\item $w^* = [1, 0]^T$
\item $d(0) = 1$
\end{itemize}

\textbf{方向2}：$\theta = \pi$（负力矩方向）
\begin{itemize}
\item 线性规划：$\max (-w_1)$
\item 解：$f = [0, 1, 1, 0]^T$
\item $w^* = [-1, 0]^T$
\item $d(\pi) = 1$
\end{itemize}

\textbf{方向3}：$\theta = \pi/2$（正力方向）
\begin{itemize}
\item 线性规划：$\max w_2$
\item 解：$f = [0, 1, 0, 1]^T$
\item $w^* = [0, 2]^T$
\item $d(\pi/2) = 2$
\end{itemize}

\textbf{方向4}：$\theta = -\pi/2$（负力方向）
\begin{itemize}
\item 线性规划：$\max (-w_2)$
\item 解：$f = [1, 0, 1, 0]^T$
\item $w^* = [0, -2]^T$
\item $d(-\pi/2) = 2$
\end{itemize}

\textbf{方向5}：$\theta = \pi/4$（力矩和力相等）
\begin{itemize}
\item 线性规划：$\max (w_1 + w_2)$
\item 需要检查所有顶点组合
\item 候选：$f = [1, 0, 0, 1]^T$，$w = [1, 0]^T$，值$= 1$
\item 候选：$f = [0, 1, 0, 1]^T$，$w = [0, 2]^T$，值$= 2$
\item 候选：$f = [1, 1, 0, 0]^T$，$w = [0, 0]^T$，值$= 0$
\item 最优：需要求解LP，但可以估计$d(\pi/4) \geq 1/\sqrt{2} \approx 0.707$
\end{itemize}

\textbf{方向6}：$\theta = 3\pi/4$（负力矩，正力）
\begin{itemize}
\item 类似分析，$d(3\pi/4) \geq 1/\sqrt{2}$
\end{itemize}

\subsection{最终质量度量}

\textbf{最小距离}：

通过检查所有关键方向，我们发现：
\begin{equation}
\text{Qual}(\{F_i\}) = \min_{\theta} d(\theta) \approx 0.707
\end{equation}

\textbf{更精确的计算}：

实际上，最小距离出现在$\theta = \pi/4$或$3\pi/4$方向附近。

\textbf{精确值}：可以通过求解以下优化问题得到：
\begin{align}
\text{最大化} \quad & d \\
\text{使得} \quad & d \leq d(\theta) \text{ 对所有 } \theta
\end{align}

\textbf{近似结果}：$\text{Qual}(\{F_i\}) \approx 0.7$

\textbf{物理解释}：
\begin{itemize}
\item 抓取可以抵抗任何方向、大小不超过0.7（归一化单位）的干扰wrench
\item 在纯力矩方向（$\theta = 0, \pi$）上，距离为1
\item 在纯力方向（$\theta = \pi/2, -\pi/2$）上，距离为2
\item 在混合方向（$\theta = \pi/4, 3\pi/4$）上，距离最小，约为0.7
\item 因此，质量度量由\textbf{最弱方向}决定
\end{itemize}

\subsection{改进抓取质量}

\textbf{问题}：当前抓取在混合方向（力矩+力）上较弱

\textbf{改进方法1}：增加接触点之间的距离
\begin{itemize}
\item 将接触点移到更远的位置
\item 例如：从$(1, 0.5)$移到$(1.5, 0.75)$
\item 这会增加力矩分量，扩展$\mathcal{C}_F$
\end{itemize}

\textbf{改进方法2}：改变接触法向量方向
\begin{itemize}
\item 使用更倾斜的法向量
\item 可以同时产生力矩和力
\item 使$\mathcal{C}_F$更"圆"，减少方向依赖性
\end{itemize}

\textbf{改进方法3}：增加接触数量
\begin{itemize}
\item 虽然4个接触已经足够形状闭合
\item 但更多接触可以提供更好的力分布
\item 提高质量度量
\end{itemize}

\section{总结}

\subsection{核心思想}

\begin{enumerate}
\item \textbf{抓取质量}需要量化，不能只判断"是否形状闭合"
\item \textbf{接触wrench集合}$\mathcal{C}_F$表示接触可以提供的所有wrench
\item \textbf{最大内接球半径}是一个简单有效的质量度量
\item \textbf{单位归一化}和\textbf{坐标系选择}是关键步骤
\end{enumerate}

\subsection{关键公式}

\begin{itemize}
\item 接触wrench集合：$\mathcal{C}_F = \left\{\sum_{i=1}^j f_i F_i \middle| f_i \in [0, f_{i,\max}] \right\}$
\item 质量度量：$\text{Qual}(\{F_i\}) = \min\{d_1, d_2, \ldots, d_k\}$
\item 单位归一化：力矩 $\to$ 力矩/$r$（特征长度）
\end{itemize}

\subsection{实际意义}

\begin{itemize}
\item 帮助选择\textbf{更好的抓取配置}
\item 评估抓取的\textbf{鲁棒性}
\item 指导\textbf{抓取规划}和\textbf{优化}
\end{itemize}

\end{document}

